%% -----------------------------------------------------------
\section{Conclusion}
%% -----------------------------------------------------------

AI mental health technologies demonstrate both remarkable therapeutic potential and serious safety risks for marginalized communities. Purpose-built interventions show that culturally-adapted, evidence-based AI can deliver meaningful benefits to underserved populations, with meta-analytic evidence confirming significant symptom reductions~\cite{abd_alrazaq_2020,li_2023}. However, systematic accessibility failures~\cite{khan_seo_2025}, pervasive demographic bias in \gls{llm}s~\cite{omar_2025,pfohl_2024}, neurotypical bias against neurodivergent users~\cite{killian_2023}, cultural misunderstanding~\cite{aleem_2024}, racial disparities in AI healthcare outcomes~\cite{haider_2024,timmons_2023}, online safety crises~\cite{tanni_2024}, harmful advice and crisis mismanagement~\cite{de_freitas_2024,he_2022,neda_2023}, and implementation gaps~\cite{smith_2023} together demonstrate that current deployment patterns deepen rather than reduce mental health inequities.

The stakes are concrete. When AI mental health tools fail, people suffer. The populations surveyed here, spanning \gls{lgbtq} youth, neurodivergent adults, blind individuals, racially and ethnically minoritized communities, culturally diverse young people, and those with limited economic resources, face elevated mental health risk, systematic barriers to traditional care, and histories of healthcare discrimination. They deserve better than being unpaid beta testers for technology deployed without adequate safety protocols. This survey points toward a path forward grounded in three principles: equity must be designed in from conception through deployment~\cite{fowler_2023,khan_seo_2025,naderbagi_2024}; safety infrastructure, including crisis protocols, human oversight, accessibility standards, and outcome monitoring, must precede deployment~\cite{de_freitas_2024,smith_2023}; and transparency, with clear information about limitations, easy access to human alternatives, and meaningful data control, must be preserved~\cite{drydakis_2022,rollwage_2024}.

The promise of AI mental health technologies for marginalized populations is real. Realizing that promise requires confronting current reality: existing approaches too often fail precisely those populations most in need. The future of equitable digital mental health depends on whether the field can move from rapid deployment to responsible development, from techno-optimism to evidence-based implementation, and from designing for imagined users to co-creating with real communities.
